\section[Notebook 1971-11: automorphisms of order $p$ trivial on $\Omega_1$]{Notebook 1971-11: automorphisms of order $p$ of an abelian $p$-group that are trivial on the elements of order $p$\protect\piste\protect\verifie{Notebook1971_11}}\label{sec:quillen-1971-11}

\emph{Context.} In 1971, Quillen studies the cohomology of finite groups and the role of their elementary abelian subgroups. \carnet{1971}{11} contains, from page~68 on, a series dated 24 August 1971, titled \q{$p$ odd}, annotated in the margin as a revision of Thompson's theorems. One reads there a lemma, then two propositions on the maximal normal elementary abelian subgroups of a $p$-group, which use it (pp.~69--71).

\emph{Notation.} For an abelian group $B$, written additively, and a prime $p$, we write $\Omega_1(B)$ for the subgroup of elements $b$ such that $pb = 0$. An \emph{abelian $p$-group} is an abelian group in which every element is killed by a power of $p$.

\emph{The page.} On page~68, Quillen states the following lemma: if $\theta$ is an automorphism of an abelian $p$-group $B$, with $\theta^p = 1$, which is the identity on $\Omega_1(B)$, then $\theta - 1$ sends $B$ into $\Omega_1(B)$. His proof, on page~69, passes to the dual group $\widehat B$, writes $\theta = 1 + N$ with $N$ having image in $p\widehat B$, and compares the terms of the binomial expansion of $(1 + N)^p = 1$; a first attempt, by a $p$-adic filtration, is crossed out. The hypothesis $p \ge 3$ serves to make the last term $N^p$ small enough. Page~69 then applies the lemma to an element of order $p$ that normalizes an abelian subgroup.

\emph{Intuition.} Put $\delta = \theta - 1$. The hypothesis says that $\delta$ kills $\Omega_1(B)$, and the conclusion that $p\delta = 0$. Since $\delta$ commutes with multiplication by $p$, it lowers the order: if $p^{k+1} b = 0$, then $p^k \delta(b) = \delta(p^k b) = 0$, since $p^k b \in \Omega_1(B)$. The powers of $\delta$ thus get smaller and smaller, and in the expansion of $(1 + \delta)^p = 1$ the term $p\delta$ can be compensated only by vanishing terms.

\begin{encadre}
\begin{theoreme}[Quillen, notebook 1971-11, p.~68]\label{thm:quillen-1971-11}
Let $p \ge 3$ be a prime, $B$ an abelian group in which every element is killed by a power of $p$, and $\theta$ an endomorphism of $B$ such that $\theta^p = 1$ and $\theta(b) = b$ for all $b \in \Omega_1(B)$. Then $p(\theta(b) - b) = 0$ for all $b \in B$, that is, $(\theta - 1)B \subset \Omega_1(B)$.
\end{theoreme}
\end{encadre}

\noindent The statement is that of the page, in a slightly more general form: $B$ is not assumed finite, and $\theta$ is an endomorphism, automatically bijective since $\theta^p = 1$.

\begin{proof}
Let $\delta = \theta - 1$. We show, by induction on $k$, that $p\delta(c) = 0$ for every $c$ such that $p^k c = 0$. For $k = 0$, $c = 0$. Suppose the result true for $k$, and let $c$ be such that $p^{k+1} c = 0$. By the lowering of the order, $p^k \delta(c) = 0$; the induction hypothesis applied to $\delta(c)$ gives $p\delta^2(c) = 0$. Then $\delta^2(c) \in \Omega_1(B)$, so $\delta^3(c) = 0$, then $\delta^m(c) = 0$ for all $m \ge 3$, and $p\delta^m(c) = 0$ for all $m \ge 2$.

Since $\theta = 1 + \delta$ and $1$ commutes with $\delta$, the binomial formula gives
\[
c = \theta^p(c) = \sum_{m = 0}^{p} \binom{p}{m}\, \delta^m(c) = c + p\,\delta(c) + \sum_{m = 2}^{p - 1} \binom{p}{m}\, \delta^m(c) + \delta^p(c).
\]
For $2 \le m \le p - 1$, the coefficient $\binom{p}{m}$ is divisible by $p$, and $p\delta^m(c) = 0$: these terms vanish. The last one, $\delta^p(c)$, vanishes because $p \ge 3$. What remains is $p\delta(c) = 0$.
\end{proof}

\noindent The proof is more direct than that of the page: it does not pass to the dual and proceeds by induction on the order of the elements. The hypothesis $p \ge 3$ serves at the same place: for $p = 2$, the last term is $\delta^2(c)$, which the induction cannot cancel.

\begin{contrexemple}[the edition's statement]\label{ce:quillen-1971-11}
For $p = 2$, the lemma is false. On $B = \mathbf Z/8$, let $\theta = -1$. Then $\theta^2 = 1$, and $\theta$ fixes $\Omega_1(B) = \{0, 4\}$, since $-4 = 4$. But $2(\theta(1) - 1) = -4 \ne 0$.
\end{contrexemple}

\noindent On $\mathbf Z/4$, on the other hand, $\theta = -1$ satisfies the conclusion ($2 \cdot (-2) = 0$): the first cyclic counterexample requires order~$8$. An exhaustive search, over the groups $\mathbf Z/9$, $\mathbf Z/27$, $\mathbf Z/9 \times \mathbf Z/3$, $\mathbf Z/9 \times \mathbf Z/9$, $\mathbf Z/27 \times \mathbf Z/3$ for $p = 3$ and $\mathbf Z/25$ for $p = 5$, found no automorphism contradicting the lemma; for $p = 2$, it finds some on $\mathbf Z/8$ and $\mathbf Z/8 \times \mathbf Z/2$, and none on $\mathbf Z/4$, $\mathbf Z/4 \times \mathbf Z/2$, $\mathbf Z/4 \times \mathbf Z/4$. The title \q{$p$ odd} of the page is therefore necessary (figure~\ref{fig:quillen-1971-11}). The hypothesis that $B$ is a $p$-group is necessary too: the ring $\mathbf Z[\zeta_p]$ of cyclotomic integers, with $\theta$ multiplication by $\zeta_p$, has $\Omega_1 = 0$ and $\theta^p = 1$, without $p(\theta - 1)$ being zero.


\begin{figure}[htbp]
\centering
\begin{tikzpicture}[>=stealth, every node/.style={font=\small}]
\begin{scope}
\foreach \k in {0,...,8} { \coordinate (n\k) at ({90-40*\k}:1.45); }
\foreach \a/\b in {1/4,4/7,7/1,2/8,8/5,5/2} { \draw[->, frise, shorten >=4pt, shorten <=4pt] (n\a) -- (n\b); }
\foreach \k in {1,2,4,5,7,8} { \fill (n\k) circle (1.7pt); \node at ({90-40*\k}:1.8) {$\k$}; }
\foreach \k in {0,3,6} { \fill[piste] (n\k) circle (2.2pt); \node[piste] at ({90-40*\k}:1.8) {$\k$}; }
\node at (0,-2.35) {$p = 3$, $B = \mathbf Z/9$, $\theta = 4$};
\end{scope}
\begin{scope}[xshift=5.6cm]
\foreach \k in {0,...,7} { \coordinate (m\k) at ({90-45*\k}:1.45); }
\foreach \k in {1,2,3} { \pgfmathtruncatemacro{\l}{8-\k} \draw[<->, frise, shorten >=4pt, shorten <=4pt] (m\k) -- (m\l); }
\foreach \k in {1,2,3,5,6,7} { \fill (m\k) circle (1.7pt); \node at ({90-45*\k}:1.8) {$\k$}; }
\foreach \k in {0,4} { \fill[piste] (m\k) circle (2.2pt); \node[piste] at ({90-45*\k}:1.8) {$\k$}; }
\node at (0,-2.35) {$p = 2$, $B = \mathbf Z/8$, $\theta = -1$};
\end{scope}
\end{tikzpicture}
\caption{The action of $\theta$ (arrows) and the subgroup $\Omega_1(B)$ (in red). On the left, an example of the edition: multiplication by $4$ on $\mathbf Z/9$ satisfies $\theta^3 = 1$, fixes $\Omega_1(B) = \{0, 3, 6\}$ and cyclically permutes $1, 4, 7$ and $2, 8, 5$; $\theta - 1$ is multiplication by $3$, with values in $\Omega_1(B)$, as the theorem requires. On the right, counterexample~\ref{ce:quillen-1971-11}: $\theta = -1$ fixes $\Omega_1(B) = \{0, 4\}$, but $\theta(1) - 1 = 6$ is not in $\Omega_1(B)$.}
\label{fig:quillen-1971-11}
\end{figure}

\emph{Verification in Lean.} The file \href{\leanblobmain/Notebook1971_11.lean}{\texttt{Notebook1971\_11.\allowbreak{}lean}} proves Theorem~\ref{thm:quillen-1971-11} under the name \texttt{nsmul\_sub\_one\_eq\_zero}, for $\theta$ in the ring \texttt{AddMonoid.End B} of endomorphisms of $B$, and under the name \texttt{addEquiv\_nsmul\_sub\_one\_eq\_zero} for an automorphism \texttt{B ≃+ B} whose $p$-th iterate is the identity, as on the page. The lowering of the order is \texttt{sub\_one\_nsmul\_eq\_zero\_of\_nsmul}, the binomial expansion evaluated at a point is \texttt{pow\_apply\_eq\_sum}. Counterexample~\ref{ce:quillen-1971-11} is \texttt{zmod8\_neg}, checked by computation. The axioms used are the three usual axioms of mathlib.

\litterature{the search did not locate a printed reference stating this lemma in this form. It is of the type of lemmas that serve Thompson's replacement theorems, which the page says it revises, and the study of elementary abelian subgroups in Quillen's 1971--1978 work on the cohomology of finite groups; the expositions of these theorems in textbooks on finite groups have not been combed. The statement is elementary, and the edition does not present it as new.}

\subsection*{The quotient of a $p$-central $p$-group by $\Omega_1$\protect\verifie{Notebook1971_11Central}}\label{sec:quillen-1971-11-a}

\emph{Notation.} For a group $P$, we write $\Omega_1(P)$ for the subgroup generated by the elements $x$ such that $x^p = 1$; it is always normal. We say that $P$ is \emph{$p$-central} if every element of order $p$ of $P$ is central, that is, if $\Omega_1(P) \subset Z(P)$. In this case, $\Omega_1(P)$ is abelian of exponent $p$ and consists exactly of the elements $x$ such that $x^p = 1$.

\emph{The page.} The series of 24 August 1971 ends, on pages~70--71 of \carnet{1971}{11}, with a proposition: for $p$ odd, if $P$ is a $p$-central $p$-group, the quotient $P/\Omega_1(P)$ is also $p$-central. Quillen's proof puts $A = \Omega_1(P)$ and chooses a subgroup $B \supset A$, normal in $P$, such that $B/A$ is a maximal normal elementary abelian subgroup of $P/A$. Raising to the $p$-th power sends $B$ into $A$; it induces an injective map from $B/A$ to $A$, compatible with conjugation since $A$ is central, from which he deduces that $B/A$ is central in $P/A$. He concludes with the preceding proposition (pp.~69--70): a maximal normal elementary abelian subgroup is a maximal elementary abelian subgroup ($p$ odd), so that $B/A$ contains all the elements of order $p$ of $P/A$. It is this intermediate proposition that uses the lemma of Theorem~\ref{thm:quillen-1971-11}.

\emph{Intuition.} An element $x$ whose class modulo $A$ has order dividing $p$ satisfies $x^p \in A$. For $g \in P$, the conjugate $u = gxg^{-1}$ has the same $p$-th power as $x$, since $x^p$ is central. If $u$ and $x$ commute modulo $A$, the group they generate has class at most~$2$, and for $p$ odd raising to the $p$-th power is a homomorphism there: $(ux^{-1})^p = u^p x^{-p} = 1$, so the commutator $gxg^{-1}x^{-1}$ is in $A$. It remains to see that $u$ and $x$ commute modulo $A$; this is the statement itself in a maximal subgroup containing $x$, whence an induction on the order.

\begin{encadre}
\begin{theoreme}[Quillen, notebook 1971-11, pp.~70--71]\label{thm:quillen-1971-11-a}
Let $p$ be an odd prime and $P$ a finite $p$-group in which every element $z$ with $z^p = 1$ is central. Then every element $y$ of $P/\Omega_1(P)$ with $y^p = 1$ is central; in other words, $\Omega_1(P/\Omega_1(P)) \subset Z(P/\Omega_1(P))$.
\end{theoreme}
\end{encadre}

\begin{proof}
Put $A = \Omega_1(P)$. Since $A = \{x \in P \mid x^p = 1\}$, we must show: if $(x^p)^p = 1$, then $(gxg^{-1}x^{-1})^p = 1$ for all $g \in P$. We argue by induction on the order of $P$; the hypothesis passes to subgroups, since an element central in $P$ is central in every subgroup containing it.

If $x$ generates $P$, the group is commutative and the commutator is trivial. Otherwise, $\langle x \rangle$ is contained in a maximal subgroup $M$, normal because $P$ is nilpotent. The conjugate $u = gxg^{-1}$ is in $M$ and satisfies $u^p = g x^p g^{-1} = x^p$, whence $(u^p)^p = 1$. The induction hypothesis, applied in $M$ to the elements $u$ and $x^{-1}$, gives $d^p = 1$ for $d = x^{-1} u x u^{-1}$; thus $d$ is central. We have $x^{-1} u = d\, u\, x^{-1}$, and an immediate induction gives, for all $n$,
\[
(u x^{-1})^n = u^n\, x^{-n}\, d^{\binom n2}.
\]
For $n = p$ odd, $p$ divides $\binom p2$, so $d^{\binom p2} = 1$ and $(ux^{-1})^p = u^p x^{-p} = 1$. Now $u x^{-1} = gxg^{-1}x^{-1}$.

Finally, if $y = xA$ satisfies $y^p = 1$, then $x^p \in A$, that is $(x^p)^p = 1$, and for all $g$ the commutator $gxg^{-1}x^{-1}$ is in $A$: $y$ commutes with $gA$.
\end{proof}

\noindent The proof keeps Quillen's first step (raising to the $p$-th power is an injective homomorphism compatible with conjugation, which forces the commutators into $A$), but replaces the appeal to the proposition of pages~69--70 by an induction on the order of $P$. It therefore does not use the lemma of Theorem~\ref{thm:quillen-1971-11}. The hypothesis $p$ odd serves at the same place as in Quillen: the divisibility of $\binom p2$ by $p$, which makes raising to the $p$-th power a homomorphism in class~$2$.

\emph{Verification in Lean.} The file \href{\leanblobmain/Notebook1971_11Central.lean}{\texttt{Notebook1971\_11Central.\allowbreak{}lean}} defines $\Omega_1(P)$ under the name \texttt{omega1}, as the subgroup generated by the elements $x$ such that $x^p = 1$, and shows that it is normal (\texttt{omega1\_normal}) and, under the hypothesis of $p$-centrality, that it consists exactly of these elements (\texttt{mem\_omega1\_iff}). The formula $(uv)^n = u^n v^n d^{\binom n2}$ is \texttt{mul\_pow\_eq\_mul\_pow\_mul\_pow\_choose}, and its case $n = p$ odd is \texttt{mul\_pow\_eq\_of\_commutator\_pow}. The induction on the order is \texttt{conj\_commutator\_pow\_eq\_one\_aux}, its statement \texttt{conj\_commutator\_pow\_eq\_one}. Theorem~\ref{thm:quillen-1971-11-a} is \texttt{pCentral\_quotient}, for elements, and \texttt{omega1\_quotient\_le\_center}, in the form of the page $\Omega_1(P) \le Z(P) \Rightarrow \Omega_1(P/\Omega_1(P)) \le Z(P/\Omega_1(P))$. The axioms used are the three usual axioms of mathlib.

Not proved here is the proposition of pages~69--70 (a maximal normal elementary abelian subgroup of a $p$-group, $p$ odd, is a maximal elementary abelian subgroup), through which Quillen's proof passes and which uses the lemma of Theorem~\ref{thm:quillen-1971-11}.

\litterature{the statement is known before the notebook: it is the theorem that J.~González-Sánchez and Th.~Weigel (2011, introduction) call Blackburn--Buckley--Thompson. It is published in this form by J.~Buckley, \q{Finite groups whose minimal subgroups are normal}, \emph{Math.~Z.} 116 (1970), 15--17, and follows from a result of Thompson that B.~Huppert reproduces as early as 1967, with a proof attributed to Blackburn (\emph{Endliche Gruppen}~I, Springer, III, Hilfssatz~12.2). Theorem~B of González-Sánchez and Weigel generalizes it to $p^i$-central pro-$p$-groups of height~$k$. The notebook reproves it in 1971, as a lemma in a revision of Thompson's theorems; only the proof by induction given here is the edition's.}
